\section{Validasi \textit{Barren Plateau}}

\subsection{\textit{Barren Plateau}}
Fenomena \textit{Barren Plateau} atau dataran tandus didefinisikan sebagai kondisi di mana lanskap fungsi biaya pada algoritma kuantum variasional menjadi sangat datar secara eksponensial seiring dengan bertambahnya jumlah \textit{qubit} dalam sistem. Kondisi patologis ini dicirikan oleh hilangnya sinyal gradien secara masif, sehingga algoritma optimasi berbasis gradien kehilangan kemampuan fungsionalnya dalam menentukan arah pergerakan menuju titik minimum energi yang diinginkan. Secara fisik, kemunculan dataran tandus ini sering kali disebabkan oleh inisialisasi parameter yang terlalu acak di dalam sirkuit kuantum yang memiliki derajat keterbelitan (\textit{entanglement}) yang sangat tinggi \cite{mcclean_2018}. Masalah skalabilitas ini menjadi hambatan utama dalam implementasi hardware era NISQ karena membatasi kedalaman sirkuit yang dapat dilatih secara efektif tanpa mengalami stagnasi proses pembelajaran. Deteksi dini terhadap kemunculan dataran tandus menjadi sangat krusial guna menjamin efisiensi komputasi pada sistem kuantum berskala besar.

Dampak dari fenomena \textit{Barren Plateau} sangat merugikan bagi efisiensi komputasi karena menuntut jumlah pengukuran (\textit{shots}) yang sangat besar hanya untuk mengekstraksi sinyal gradien yang sangat kecil dari dominasi derau hardware. Ketika varians dari gradien tersebut menyusut di bawah batas presisi alat ukur, sirkuit kuantum secara efektif akan kehilangan kapasitasnya untuk melakukan pembaruan parameter yang bermakna bagi proses konvergensi. Fenomena ini memaksa komunitas peneliti untuk mengembangkan berbagai teknik inisialisasi cerdas guna menjaga agar gradien tetap berada pada level yang dapat terdeteksi. Dalam penelitian ini, integrasi teori permainan melalui \textit{Exact Potential Game} (EPG) diusulkan sebagai solusi preventif guna memberikan inisialisasi parameter yang lebih terarah dan keluar dari zona tandus tersebut. Dengan memanfaatkan titik ekuilibrium Nash sebagai posisi awal, algoritma VQE dapat melewati fase pencarian acak yang rentan terhadap stagnasi gradien di ruang Hilbert yang luas.

\subsection{Entropi Von Neumann}
Entropi Von Neumann dimanfaatkan sebagai metrik kuantitatif yang sangat penting untuk mengukur derajat keterbelitan (\textit{entanglement}) di dalam sistem kuantum multiqubit yang sedang dioptimasi. Konsep ini diperluas dari paradigma entropi Shannon guna mengukur derajat kemurnian atau campuran suatu keadaan kuantum melalui operator densitas $\rho$ yang dinyatakan dalam bentuk $S(\rho) = -\text{Tr}(\rho \log \rho)$. Dalam arsitektur \textit{Parametrized Quantum Circuits} (PQC), peningkatan keterbelitan yang tidak terkontrol sering kali berkorelasi kuat dengan hilangnya lanskap gradien secara sistemik di seluruh ruang parameter. Secara teoretis, sirkuit yang telah mencapai karakteristik \textit{2-design} cenderung memiliki nilai entropi yang mendekati nilai maksimum atau yang dikenal sebagai \textit{Page entropy}. Nilai entropi yang sangat tinggi ini menandakan bahwa informasi kuantum telah tersebar secara merata di seluruh ruang Hilbert, sehingga menyebabkan fungsi biaya menjadi sulit untuk dieksplorasi oleh algoritma optimasi klasik.

Nilai numerik dari entropi Von Neumann memberikan indikasi langsung mengenai status fisik sistem, di mana nilai $S=1$ mengonfirmasi terjadinya keterbelitan maksimal yang merepresentasikan Keadaan Bell. Keberadaan nilai entropi di rentang $0 < S < 1$ menunjukkan kondisi keterbelitan parsial (\textit{partial entanglement}), yang menandakan adanya korelasi kuantum antar-\textit{qubit} namun belum mencapai titik jenuh informasi sistemik. Sebaliknya, nilai $S=0$ menunjukkan bahwa sistem berada dalam keadaan terpisahkan (\textit{separable state}) tanpa adanya korelasi kuantum antar-sub-sistem yang dapat diobservasi. Pemanfaatan entropi sebagai alat diagnostik memungkinkan peneliti untuk memantau apakah sirkuit sedang mendekati ambang batas \textit{2-design} yang memicu munculnya fenomena \textit{Barren Plateau}. Dengan demikian, entropi berfungsi sebagai parameter kontrol untuk menjaga agar sistem tetap berada dalam zona yang dapat dioptimasi secara efektif guna mencapai konvergensi solusi yang diinginkan.

Dinamika nilai entropi Von Neumann selama proses iterasi optimasi mencerminkan evolusi struktural pada vektor keadaan sistem seiring dengan pembaruan parameter yang dilakukan oleh optimizer. Penurunan nilai entropi dari kondisi maksimal ($S \approx 1$) menuju titik deterministik ($S \to 0$) menandakan transisi sistem dari status keterbelitan acak menuju konvergensi solusi portofolio yang bersifat terpisahkan. Sebaliknya, peningkatan entropi selama tahapan iterasi menunjukkan bahwa sistem sedang mengeksplorasi ruang keterbelitan guna meminimalkan fungsi biaya, yang berakibat pada distribusi probabilitas yang tersebar secara merata. Fenomena ini secara empiris membuktikan bahwa derajat keterbelitan di dalam sirkuit bersifat dinamis dan dapat bertransformasi sebagai respons terhadap pencarian titik minimum pada lanskap energi portofolio. Observasi terhadap fluktuasi entropi ini memberikan wawasan mendalam mengenai mekanisme internal sirkuit kuantum dalam memproses informasi korelasi aset finansial selama siklus optimasi hibrida berlangsung.

\subsection{Varians Gradien}
Fenomena \textit{Barren Plateau} secara formal dapat didefinisikan melalui analisis terhadap perilaku statistik dari gradien fungsi biaya terhadap parameter-parameter sirkuit yang dioptimasi secara iteratif. Pada kondisi ini, nilai rata-rata dari gradien tersebut cenderung mendekati nol, namun parameter yang paling krusial adalah variansnya yang menghilang secara eksponensial terhadap penambahan jumlah \textit{qubit}. Secara matematis, varians gradien $\text{Var}[\partial_{\theta} E]$ akan menyusut menuju nol mengikuti pola $\mathcal{O}(2^{-n})$, di mana variabel $n$ merepresentasikan jumlah \textit{qubit} yang digunakan dalam arsitektur komputer kuantum. Penyusutan varians yang bersifat eksponensial ini mengakibatkan sinyal gradien terkubur di bawah level derau perangkat (\textit{hardware noise}) serta batasan presisi numerik yang tersedia. Oleh karena itu, mitigasi terhadap penyusutan varians gradien menjadi prasyarat mutlak bagi keberhasilan implementasi algoritma kuantum pada masalah-masalah finansial yang melibatkan jumlah aset yang signifikan.
